% Research proposal: the structure most fellowship and grant calls ask for.
\documentclass[11pt,a4paper]{article}

\usepackage[T1]{fontenc}
\usepackage[margin=2.5cm]{geometry}
\usepackage{amsmath}
\usepackage{graphicx}
\usepackage{booktabs}
\usepackage{enumitem}
\usepackage[hidelinks]{hyperref}

\title{Proposal Title}
\author{Applicant Name \\ Proposed supervisor: Supervisor Name \\ Department, University}
\date{\today}

\begin{document}

\maketitle

\section*{Summary}
Four or five sentences a reviewer outside the subfield can follow: the gap,
the idea, the plan and the expected outcome.

\section{Background and Motivation}
What is known, what is missing and why it matters now~\cite{lamport1994}.

\section{Research Questions}
\begin{enumerate}[label=\textbf{Q\arabic*.}]
  \item First question, stated so that an experiment or proof can answer it.
  \item Second question.
\end{enumerate}

\section{Methodology}
\subsection{Approach}
How each question will be answered: data, instruments, models.

\subsection{Preliminary Results}
Evidence that the approach is feasible.

\section{Work Plan}

\begin{table}[htbp]
  \centering
  \caption{Timeline by quarter.}
  \begin{tabular}{llp{7.5cm}}
    \toprule
    Period & Phase & Milestone \\
    \midrule
    Q1--Q2 & Setup       & Literature review complete; apparatus commissioned \\
    Q3--Q4 & Experiments & First data set; conference abstract \\
    Q5--Q6 & Analysis    & Model validated; journal manuscript submitted \\
    \bottomrule
  \end{tabular}
\end{table}

\section{Expected Outcomes and Impact}
Publications, data sets, software, and who benefits.

\section{Resources and Budget}
Equipment, travel and consumables, each with a one-line justification.

\bibliographystyle{plain}
\bibliography{references}

\end{document}
